\chapter{Why We Do This: Writing for Everyone}
\label{chap:why}

\begin{quote}
\enquote{Literature is my Utopia. Here I am not disenfranchised. No
barrier of the senses shuts me out from the sweet, gracious discourse
of my book friends.}
--- Helen Keller, \emph{The Story of My Life}~\autocite{keller1903}
\end{quote}

Years from now, a student may open your thesis at two in the morning.
You may never meet that reader. One reader may see the page. Another may listen through a screen reader. Someone else may follow the text with her fingertips on a braille display.

Whether your words can reach them is decided now. It may cost an
afternoon, a good coffee, and a few choices made while you write. This
chapter is about why that afternoon matters. It begins in Alabama,
beside a water pump.

%=====================================================================
\section{A word at the pump}
\label{sec:keller}

Helen Keller was born in 1880 in Tuscumbia, Alabama. When she was
nineteen months old, an illness took away her sight and her hearing.
For the next five years, she later wrote, she lived \enquote{at sea in
a dense fog}~\autocite{keller1903}.

Then came an April morning in 1887. Anne Sullivan took her to the
family's water pump. Water ran over one of Helen's hands. Into the
other, Sullivan spelled \mbox{w-a-t-e-r}. Again and again.

Something opened.

\enquote{Everything had a name,} Keller later remembered, \enquote{and
each name gave birth to a new thought.} The child did not become
intelligent on that morning. She had always been intelligent. What changed was the form in which the world reached her.

Keller went on to Radcliffe College, where in 1904 she became the first
deafblind person to earn a Bachelor of Arts degree. She wrote more than
a dozen books. But before the books, there was a word. Before the word,
there was a hand waiting to understand.

\begin{figure}[htbp]
\centering
\includegraphics[width=0.52\linewidth,
alt={Helen Keller, around age forty, seated in a wooden chair in front
of a shingled wall, cradling a large magnolia blossom in both hands.
She wears a dark embroidered dress and a pendant necklace and faces
slightly to the side, calm and composed.}]%
{figures/helen-keller.jpg}
\caption[Helen Keller]{Helen Keller (1880--1968), author and advocate,
born in Tuscumbia, Alabama. (Photograph: Los
Angeles Times, c.~1920.)}
\label{fig:keller}
\end{figure}

%=====================================================================
\section{A mother's voice}
\label{sec:pontryagin}

Lev Semyonovich Pontryagin was born in Moscow in 1908. At fourteen, a
stove explosion took his sight. His mother, Tatyana Andreyevna, was a
seamstress. She had no mathematical training. But she sat beside him
and became his reader.

For years she read mathematical papers aloud to him. When she met an
unfamiliar symbol, she described its shape. When the printed notation
had no natural sound, they gave it a name of their own~\autocite{mactutorpontryagin}.

Pontryagin duality, Pontryagin classes, the Pontryagin maximum
principle---these now belong to mathematics. But before they entered
textbooks and lecture halls, they passed through a mother's voice.

Access, for Pontryagin, took the form of someone sitting beside him
and telling him, line by line, what was on the page.

\begin{figure}[htbp]
\centering
\includegraphics[width=0.6\linewidth,
alt={Black-and-white portrait of Lev Pontryagin, silver-haired, in a
tweed jacket and dark tie, seated at a desk with his hands clasped
over papers, his head turned in near profile.}]%
{figures/pontryagin.jpg}
\caption[Lev Pontryagin]{Lev Semyonovich Pontryagin (1908--1988), who
lost his sight at fourteen. (Photograph: unknown author, Moscow
State University history of mathematics collection; CC~BY~4.0, via
Wikimedia Commons.)}
\label{fig:pontryagin}
\end{figure}

%=====================================================================
\section{When the page went dark}
\label{sec:euler}

Leonhard Euler lost the sight of his right eye around 1738. He is said
to have remarked, \enquote{Now I will have fewer
distractions}~\autocite{calinger2016}. Later, a cataract in his left
eye took most of the rest. From 1766 onward he was almost totally
blind.

Still, the work continued. Nearly half of his roughly eight hundred and
fifty writings were produced after 1766, dictated to assistants. He
calculated in memory. He composed without seeing the page. The formulas
came from his mind, but they reached the world through other people's
hands.

\begin{figure}[htbp]
\centering
\includegraphics[width=0.5\linewidth,
alt={Pastel portrait of Leonhard Euler in a grey cap and a blue
dressing gown over a white ruffled shirt, against a dark background.
His left eye is bright and alert; his right eyelid droops nearly
closed.}]%
{figures/euler.jpg}
\caption[Leonhard Euler]{Leonhard Euler (1707--1783), in the 1753
pastel by Jakob Emanuel Handmann. (Kunstmuseum Basel; public domain,
via Wikimedia Commons.)}
\label{fig:euler}
\end{figure}

%=====================================================================
\section{A way in}
\label{sec:bridge}

Put the three stories beside one another. Sullivan spelled into
Keller's hand. Tatyana Pontryagin read mathematics she did not fully understand. Euler's assistants wrote down work that he could no
longer see.

In each story, knowledge had to change form before it could travel. For Keller, language entered through touch. Pontryagin received the
printed page through his mother's voice. Euler's thoughts reached the
page through dictation.

That is easy to miss when we talk about accessibility today, because
the helper is often no longer a person in the room. Today, the intermediary may be a screen reader such as VoiceOver, a
braille display, or software that reads a document aloud. But the old problem remains: the
reader still needs the document to arrive in a form they can use.

Software cannot guess what the author meant. It reads what the file
contains. If a figure has no description, a screen reader may say only
\enquote{image} or \enquote{Figure 1.2}. That tells the reader that
something is on the page, but not what the figure shows, why it was
included, or how it supports the argument.

This is why structure matters. Alt text tells the reader what the figure shows and why it matters.
 A table header tells the software which entries belong together. Readable mathematics gives a
formula a chance to be spoken, searched, or sent to braille.

The point is not to make the PDF look more polished. The point is to
make the document usable when sight is not the only way in.

%=====================================================================
\section{An afternoon}
\label{sec}

The template supplies much of the structure. The author still has to
describe the figures, use headings correctly, and give tables the
information that software needs. Over a dissertation, this may add up
to an afternoon, perhaps two.

A dissertation is one of the most permanent documents most of us will
write. It may remain in a university repository for decades. A future
reader will not know how much time the author spent preparing the file.
The reader will know only whether the document works.

Making it accessible is part of making that document complete. The
Auburn Creed gives language to the responsibility behind this work:
\enquote{I believe in the human touch, which cultivates sympathy with
my fellow men and mutual helpfulness and brings happiness for
all}~\autocite{auburncreed}.

The next chapter turns from why this matters to how it is done.
